\apartado{2.5}{Reglas básicas}{calculo}
\bajada{Cuándo se suman las probabilidades y cuándo se multiplican. Y por qué
suponer independencia sin verificarla produce números que parecen razonables y
no lo son — algo que en psicología pasa todo el tiempo, porque casi nada es
independiente.}

\definicion{Mutuamente excluyentes e independientes}{%
  \textbf{Mutuamente excluyentes}: no pueden ocurrir a la vez. Si ocurre uno,
  el otro queda descartado.\par
  \textbf{Independientes}: saber que ocurrió uno no cambia la probabilidad del
  otro.\par
  No son lo mismo, y casi son opuestos: si dos eventos son excluyentes, saber
  que ocurrió uno te dice \emph{con certeza} que el otro no ocurrió, o sea que
  se influyen al máximo.}

\minihistoria{La comorbilidad es el ejemplo, no una excepción}{%
  La depresión y la ansiedad aparecen juntas mucho más seguido de lo que
  ocurriría por azar. Lo mismo pasa con consumo y trastornos del ánimo, o con
  ansiedad e insomnio.\par
  Por eso multiplicar probabilidades «porque sí» es un error caro en
  psicología: da números más chicos que la realidad justo en los casos más
  graves, que son los que más recursos necesitan.}

\ejercicio{0}{Explica con tus palabras qué significa que dos eventos sean
  \emph{mutuamente excluyentes}, y da un ejemplo de una historia clínica.
  Después decide: ¿«presentar sintomatología depresiva» y «presentar
  sintomatología ansiosa» son mutuamente excluyentes?}{}
\renglones[4]
\respuesta{Que no pueden ocurrir a la vez. No son excluyentes: una misma
persona puede presentar las dos.}
\solucion{%
  Mutuamente excluyentes significa que \textbf{no pueden ocurrir al mismo
  tiempo}: formalmente, $P(A\cap B)=0$.
  \par
  Un ejemplo de historia clínica: el diagnóstico principal registrado en la
  ficha. Sólo puede ser uno, así que «diagnóstico principal: depresión mayor» y
  «diagnóstico principal: trastorno de pánico» sí son excluyentes.
  \par
  Pero \textbf{presentar sintomatología depresiva y ansiosa no lo es}: una
  misma persona puede presentar las dos, y de hecho es lo más frecuente. Ése es
  el corazón del apartado.}
\enclase{Que alguien nombre un par de eventos realmente excluyentes en
psicología. Cuesta más de lo que parece, y esa dificultad ya enseña algo.}

\ejercicio{1}{En un servicio se registran dos síntomas. De 200 fichas, 43
  presentan sintomatología depresiva y 21 ansiosa.
  \textbf{Antes de calcular nada, escribe cuántas fichas crees que presentan
  las dos:} \underline{\hspace{2cm}}
  \par\vspace{0.2em}
  Guarda ese número: lo vas a comparar en el ejercicio 6.}{Predice primero}
\renglones[3]
\respuesta{Son 17. La mayoría estima bastante menos.}
\solucion{%
  La respuesta es \textbf{17}, y se comprueba en el ejercicio 6.
  \par
  Casi todo el mundo estima entre 4 y 8, porque razona «si son 43 y 21 de 200,
  coincidir debe ser raro». Ese razonamiento supone independencia sin decirlo,
  y bajo independencia darían unas 4,5 fichas.
  \par
  La distancia entre lo que estimaste y el 17 es exactamente el tamaño de la
  comorbilidad. No es un detalle estadístico: es un hecho clínico.}
\enclase{Anota dos o tres estimaciones en el pizarrón. Vuelve a ellas en el
ejercicio 6 — es donde se produce el aprendizaje.}

\ejercicio{1}{Sea $A$ = «presenta sintomatología ansiosa» y $D$ = «presenta
  sintomatología depresiva», sobre las 200 fichas, con $|A|=21$, $|D|=43$ y
  $|A\cap D|=17$. Calcula $P(A)$, $P(D)$ y $P(A\cap D)$.}{}
\renglones[3]
\respuesta{$P(A)=0{,}105$; $P(D)=0{,}215$; $P(A\cap D)=0{,}085$.}
\solucion{%
  \[ P(A)=\frac{21}{200}=0{,}105 \qquad P(D)=\frac{43}{200}=0{,}215
     \qquad P(A\cap D)=\frac{17}{200}=0{,}085 \]
  La intersección son las personas que presentan \emph{las dos} cosas, no las
  que presentan alguna.}

\ejercicio{2}{Dos preguntas que suenan igual y no lo son. Contesta las dos y
  explica qué cambia.
  \begin{enumerate}[label=(\alph*),topsep=2pt,itemsep=1pt]
    \item ¿Qué probabilidad hay de que una ficha presente \textbf{las dos}
      sintomatologías?
    \item ¿Qué probabilidad hay de que presente \textbf{al menos una}?
  \end{enumerate}}{Dos casos casi iguales}
\renglones[5]
\respuesta{(a) $17/200 = 0{,}085$. (b) $47/200 = 0{,}235$.}
\solucion{%
  (a) Es la intersección: $\dfrac{17}{200}=0{,}085$.
  \par
  (b) Es la unión, y hay que restar la intersección:
  \[ P(A\cup D)=P(A)+P(D)-P(A\cap D)
     =\frac{21}{200}+\frac{43}{200}-\frac{17}{200}=\frac{47}{200}=0{,}235 \]
  \par
  En personas: $21+43-17=47$. \textbf{Sin restar los 17} darías 64,
  contándolos dos veces — un 36\,\% de más.
  \par
  «Las dos» y «al menos una» son preguntas clínicas distintas: la primera
  identifica los casos complejos, la segunda dimensiona la demanda total del
  servicio.}

\ejercicio{2}{Un informe de un servicio dice:
  \begin{quote}\itshape
  «El 10,5\,\% presenta sintomatología ansiosa y el 21,5\,\% depresiva. Como
  son independientes, el 32\,\% presenta al menos una de las dos.»
  \end{quote}
  Hay \textbf{dos} errores en esa frase. Señala los dos y explica cada uno.}{Encuentra el error}
\renglones[6]
\respuesta{(1) Suma sin restar la intersección. (2) Afirma independencia sin
verificarla, y no son independientes. El valor real es 23,5\,\%.}
\solucion{%
  \textbf{Primer error}: sumó 10,5 más 21,5 y obtuvo 32 sin restar a quienes
  están en los dos grupos. Aunque fueran independientes, la regla de la suma
  exige restar la intersección.
  \par
  \textbf{Segundo error}: afirma la independencia en vez de verificarla. Y no
  son independientes — se comprueba en el ejercicio siguiente.
  \par
  El valor correcto es $\dfrac{47}{200}$, o sea 23,5\,\%, bastante menos que el
  32\,\% del informe.
  \par
  \textit{Lo llamativo}: los dos errores empujan en la misma dirección
  —sobreestimar— así que ninguno corrige al otro. Un servicio que planifique
  con ese 32\,\% pide más recursos de los que necesita para la demanda total, y
  al mismo tiempo subestima los casos complejos, como muestra el ejercicio 8.}

\ejercicio{3}{¿Son \textbf{independientes} presentar sintomatología depresiva y
  presentar sintomatología ansiosa?
  \begin{enumerate}[label=(\alph*),topsep=2pt,itemsep=1pt]
    \item ¿Cuántos casos con ambas \emph{esperarías} si fueran independientes?
    \item ¿Cuántos hay en realidad?
    \item Calcula $P(A\mid D)$ y compárala con $P(A)$.
    \item Compara con tu estimación del ejercicio 2.
  \end{enumerate}}{}
\renglones[7]
\respuesta{(a) Unas 4,5. (b) 17. (c) $17/43\approx 0{,}395$ contra $0{,}105$.
No son independientes.}
\solucion{%
  (a) Bajo independencia: $0{,}105\times 0{,}215\times 200 \approx 4{,}5$
  fichas.
  \par
  (b) En realidad hay \textbf{17}: casi cuatro veces más.
  \par
  (c) $P(A\mid D)=\dfrac{17}{43}\approx 0{,}395$, contra $P(A)=0{,}105$.
  \par
  \textbf{Saber que alguien presenta sintomatología depresiva casi cuadruplica
  la probabilidad de que presente también ansiosa.} Eso es exactamente lo
  contrario de la independencia, y es un hallazgo clínico real: la comorbilidad
  entre depresión y ansiedad está documentada en toda la literatura.
  \par
  (d) Si estimaste entre 4 y 8, estabas razonando como si fueran
  independientes sin darte cuenta. Es el sesgo por defecto, y por eso conviene
  verificar siempre.}
\enclase{Éste es el momento del apartado. Vuelve a las estimaciones del
pizarrón: la mayoría habrá dicho un número cercano al 4,5 de la independencia,
sin haber hecho ninguna cuenta.}

\ejercicio{3}{Un investigador quiere calcular la probabilidad de que un
  paciente con sintomatología depresiva \emph{desarrolle} sintomatología
  ansiosa en los seis meses siguientes. Tiene las 200 fichas con los dos
  puntajes tomados el mismo día.
  \textbf{¿Alcanzan los datos?} Si te parece que sí, calcula; si no, escribe
  qué falta.}{¿Alcanzan los datos?}
\renglones[5]
\respuesta{No alcanzan: las fichas son de un solo momento y la pregunta es
sobre lo que ocurre después. Haría falta una segunda medición.}
\solucion{%
  \textbf{No alcanzan.} Las dos medidas se tomaron \emph{el mismo día}, así que
  el archivo describe lo que pasa a la vez, no lo que pasa después.
  \par
  El $17/43 \approx 0{,}395$ del ejercicio anterior es tentador y responde otra
  pregunta: «entre los que hoy tienen síntomas depresivos, qué proporción hoy
  también tiene ansiosos». No dice nada sobre desarrollar algo en el futuro.
  \par
  Haría falta un \textbf{seguimiento}: volver a medir a los seis meses,
  distinguiendo a quienes ya tenían síntomas ansiosos al inicio. Confundir un
  dato transversal con uno longitudinal es uno de los errores más frecuentes en
  la lectura de investigación en psicología.}

\ejercicio{4}{Un equipo planifica el servicio del año que viene y escribe:
  \begin{quote}\itshape
  «Esperamos unos 4 o 5 casos con comorbilidad, que son los que requieren
  seguimiento intensivo. Reservamos dos horas semanales para eso.»
  \end{quote}
  \begin{enumerate}[label=(\alph*),topsep=2pt,itemsep=1pt]
    \item ¿De dónde salió ese «4 o 5», y qué supuso quien lo calculó?
    \item ¿Cuál es el número real?
    \item ¿Qué consecuencia concreta tiene el error para el equipo?
  \end{enumerate}}{}
\renglones[7]
\respuesta{(a) De multiplicar las probabilidades, suponiendo independencia.
(b) 17. (c) Respuesta abierta: falta casi cuatro veces el tiempo previsto.}
\solucion{%
  (a) De $0{,}105\times 0{,}215\times 200 \approx 4{,}5$. Quien lo calculó
  \textbf{supuso independencia} sin verificarla — probablemente sin notar que
  lo estaba suponiendo.
  \par
  (b) El número real es \textbf{17}.
  \par
  (c) El equipo reservó dos horas semanales para unos 4 casos y va a recibir
  17. Y no son casos cualesquiera: la comorbilidad suele implicar cuadros más
  complejos, más sesiones y mayor riesgo. El error no deja al servicio un poco
  corto — lo deja sin recursos justo donde más hacen falta.
  \par
  \textbf{La regla}: multiplicar probabilidades sólo se puede si son
  independientes, y la independencia se verifica, no se supone. En psicología
  casi nada es independiente, así que la carga de la prueba está siempre del
  lado de quien multiplica.}

\trampa{multiplicar dos probabilidades para obtener la de que ocurran las dos}%
  {es la fórmula que se aprende primero y funciona con dados y monedas, que es
   donde todos la practicamos}%
  {sólo vale bajo independencia. Con dos medidas tomadas sobre las mismas
   personas, lo esperable es que NO sean independientes: compruébalo antes de
   multiplicar}

\cierre{%
  La regla de la suma pide restar la intersección y la del producto pide
  verificar la independencia. Las dos se olvidan por la misma razón: el
  resultado que sale al olvidarlas parece razonable.
  \par
  Y quedó planteada una pregunta que este apartado no puede responder. Sabemos
  que $P(A\mid D)$ es casi cuatro veces $P(A)$ — pero cuando lo que conocemos
  es el resultado de un test y lo que queremos saber es la condición, la
  condicional hay que darla vuelta. Ése es el apartado siguiente.}
